\begin{textsample}{NEG Dim 5 – human – Score -9.00 – mg/mg\_ef\_linguagens\_arte\_5\_p7.txt}  \label{ex:f5_neg_017}
Este currículo também foi pensado nessas \textbf{especificidades} e particularidades regionais, e não alterou a sugestão da BNCC que possibilita uma maleabilidade dos \textbf{objetos} de \textbf{conhecimento} a serem \textbf{trabalhados} de acordo com os \textbf{anos} de escolaridade.
O que se sugere, contudo, é que em cada uma das linguagens haja um aprofundamento gradativo da \textbf{complexidade} das \textbf{habilidades} a serem \textbf{trabalhadas}, à medida em que o \textbf{estudante} avance nos \textbf{anos} de escolaridade, respeitando a área de formação do docente e seu \textbf{desenvolvimento} pedagógico.

% matched lemmas: ano, complexidade, conhecimento, desenvolvimento, especificidade, estudante, habilidade, objeto, trabalhar
\end{textsample}
